% formalisation.tex: blueprint extension of an accepted paper (AGENTS.md, section 6.6, case c).
% Not refereed: Lean checks every node. paper.tex is never edited.
% Nodes of paper.tex are referred to by their labels: xr-hyper reads paper.aux, so compile
% paper.tex first (CI does). New labels start with f (flem:, fprop:, fdef:) and are never reused.
% Every node has \lean{...} with its Lean name and an exhaustive \uses{...}; never write \leanok.
% Compile: latexmk -pdf -interaction=nonstopmode -halt-on-error -cd papers/<P>/formalisation.tex
\documentclass[11pt]{article}
\usepackage[T1]{fontenc}
\usepackage[utf8]{inputenc}
\usepackage{lmodern}
\usepackage[a4paper,margin=2.5cm]{geometry}
\usepackage{amsmath,amssymb,mathtools}
\usepackage{xcolor}
\usepackage{xr-hyper}
\usepackage[colorlinks=true,linkcolor=blue!50!black,citecolor=blue!50!black,urlcolor=blue!50!black]{hyperref}
\usepackage{../../templates/paper/blueprint-macros}
\externaldocument{paper}
\counterwithout{theorem}{section}
\renewcommand{\thetheorem}{F\arabic{theorem}}

% The paper's macros, copied from paper.tex.
\newcommand{\E}{\mathbb{E}}

\title{Formalisation notes for P003-transfer-setting}
\author{Incognita agents\thanks{Not refereed: every node is checked in Lean. Numbers F1, F2, \dots{}
  are this file's; other numbers refer to the paper.}}
\date{2026-11-17}

\begin{document}
\maketitle

\section{Extra lemmas}

\begin{lemma}[The prior game, round by round]\label{flem:prior-state}
\lean{Incognita.P003.measurable_state, Incognita.P003.sum_state_fst,
  Incognita.P003.state_fst_le_succ}
\uses{def:prior-game, def:source-space}
For a prior strategy and every $r\in\mathbb N$, the state after $r$ rounds (the counts $C(r)$ and
the history $(A'_i,Y'_i)_{i<r}$) is a measurable function of $\omega\in\Omega'$; the counts satisfy
$\sum_k C_k(r)=r$ and $C_k(r)\le C_k(r+1)$ for every arm $k$.
\end{lemma}

\begin{proof}
\uses{def:prior-game, def:source-space}
Induction on $r$. The arm $A'_r$ is a measurable function of $U'$ and of the history; the new counts
are a function of $(A'_r, C(r))$, with values in the countable set $[K]\times\mathbb N^K$, hence
measurable; the observation $X'_{A'_r,C_{A'_r}(r)}$ is measurable as a function of the stacks and
of $(A'_r,C(r))$ (a countable family of measurable sections). Only the count of $A'_r$ increases,
by $1$.
\end{proof}

\begin{lemma}[The game reads only the box of its counts]\label{flem:box-claim}
\lean{Incognita.P003.state_eq_of_eqOn_box, Incognita.P003.stoppedCounts_eq_of_eqOn_box}
\uses{def:prior-game, def:stopped-game, def:source-space, flem:prior-state}
Let $\omega,\omega'\in\Omega'$ have the same randomisation and the same values $X'_{k,j}$ for
$(k,j)\in B(n)$. If $C(r)\le n$ in the game run on $\omega$, the first $r$ rounds of the games run on
$\omega$ and on $\omega'$ coincide. For a stopping time $\tau$ of the prior game, if the stopped
counts at $\omega$ equal $n$, so do the stopped counts at $\omega'$.
\end{lemma}

\begin{proof}
\uses{def:prior-game, def:stopped-game, flem:prior-state}
This is the claim of the proofs of Lemma~\ref{lem:prior-game} and
Proposition~\ref{prop:stopped-games}~(a). The first part by induction on $r$: at round $r$ the two
games choose the same arm $a$ with the same count $C_a(r)$, and $C_a(r)+1=C_a(r+1)\le n_a$, so the
value read is in the box. For the second, with $r_0=\tau(\omega)$, the first $r_0$ rounds coincide;
$\{\tau=r_0\}\in\mathcal F'_{r_0}$ is a function of the randomisation and of these rounds, so
$\tau(\omega')=r_0$ and the counts at $r_0$ agree.
\end{proof}

\begin{lemma}[Stopped counts form a size rule]\label{flem:stopped-size-rule}
\lean{Incognita.P003.stoppedCounts_isSizeRule}
\uses{def:size-rule, def:stopped-game, flem:prior-state, flem:box-claim}
Item (a) of Proposition~\ref{prop:stopped-games}, proved without Lemma~\ref{lem:prior-game}.
\end{lemma}

\begin{proof}
\uses{def:size-rule, def:source-space, def:stopped-game, flem:prior-state, flem:box-claim}
A stopping time is measurable, since $\mathcal F'_r$ is contained in the $\sigma$-algebra of
$\Omega'$ by Lemma~\ref{flem:prior-state}; so the stopped counts
$\omega\mapsto C(\tau(\omega))(\omega)$ are measurable (a countable family of measurable sections,
indexed by $r$), and their total is $\tau\le T'$. For $n\in\mathbb N^K$ let $\hat\omega$ keep the
randomisation and the values of $B(n)$ of $\omega$ and put $0$ elsewhere; $\omega\mapsto\hat\omega$ is
$\mathcal G_n$-measurable. By Lemma~\ref{flem:box-claim} applied to $(\omega,\hat\omega)$ and to
$(\hat\omega,\omega)$, the event $\{N'=n\}$ is its own preimage under this map, hence in
$\mathcal G_n$.
\end{proof}

\section{New decompositions of the paper's proofs}

\begin{proof}[Formal route to Lemma~\ref{lem:prior-game}]
\proves{lem:prior-game}
\uses{def:size-rule, def:prior-game, def:stopped-game, flem:stopped-size-rule}
(a) as in the paper. (b) is Lemma~\ref{flem:stopped-size-rule} for the constant stopping time
$\tau\equiv T'$ (\texttt{PriorStrategy.isStoppingTime\_const}): the counts are a size rule with total
at most $T'$, so each count is at most $T'$.
\end{proof}

\end{document}
